% This Source Code Form is subject to the terms of the Mozilla Public
% License, v. 2.0. If a copy of the MPL was not distributed with this
% file, You can obtain one at https://mozilla.org/MPL/2.0/.

\documentclass[a4paper,12pt]{article}
\usepackage[utf8]{inputenc}
\usepackage[T2A]{fontenc}
\usepackage[english,russian]{babel}
\usepackage{geometry}
\geometry{left=2cm,right=2cm,top=2cm,bottom=2cm}
\usepackage{amsmath}
\usepackage{multirow}
\usepackage{array}
\usepackage{float}

\begin{document}
	
	\subsubsection*{Опыт 1. Взаимодействие металлов с хлороводородной (соляной) и разбавленной серной кислотами}
	
	\subsubsection*{Ход работы}
	
	\quad В четыре пробирки помещены порознь стружки алюминия, железа, свинца и меди. В каждую пробирку налито приблизительно 2 мл разбавленной соляной кислоты (\(C = 2 \, \text{моль/л}\)). Пробирки слегка подогреты. Наблюдения за выделением водорода:
	
	\begin{itemize}
		\item Алюминий: активное выделение газа.
		\item Железо: умеренное выделение газа.
		\item Свинец: очень слабое выделение газа (практически не наблюдается).
		\item Медь: выделение газа отсутствует.
	\end{itemize}
	
	\quad Аналогичный опыт проведен с разбавленной серной кислотой (\(C = 1 \, \text{моль/л}\)). Результаты аналогичны.
	
	\subsubsection*{Уравнения реакций}
	\begin{align*}
		& 2\text{Al} + 6\text{HCl} \rightarrow 2\text{AlCl}_3 + 3\text{H}_2 \uparrow \\
		& \text{Fe} + 2\text{HCl} \rightarrow \text{FeCl}_2 + \text{H}_2 \uparrow \\
		& 2\text{Al} + 3\text{H}_2\text{SO}_4 \rightarrow \text{Al}_2(\text{SO}_4)_3 + 3\text{H}_2 \uparrow \\
		& \text{Fe} + \text{H}_2\text{SO}_4 \rightarrow \text{FeSO}_4 + \text{H}_2 \uparrow
	\end{align*}
	
	\subsubsection*{Объяснение}
	
	\quad Свинец и медь практически не реагируют с разбавленными кислотами (HCl, \(H_2SO_4\)) из-за низкой активности (\(E^0_{\text{Pb}^{2+}/\text{Pb}} > 0\), \(E^0_{\text{Cu}^{2+}/\text{Cu}} > 0\)) и пассивирующего действия образующихся хлоридных и сульфатных слоев.
	
	\subsubsection*{Опыт 2. Взаимодействие меди и алюминия с концентрированной серной кислотой}
	
	\subsubsection*{Ход работы}
	
	\quad В пробирку налито около 1 мл концентрированной \(H_2SO_4\) (\(\omega = 95.5\%\)). Опущены 1–2 стружки меди, пробирка осторожно нагрета. Наблюдается выделение газа с резким запахом \(SO_2\). Аналогичный опыт с алюминием проведен с осторожностью: выделение газа может сопровождаться изменением запаха (возможно выделение \(H_2S\)).
	
	\subsubsection*{Уравнения реакций}
	\begin{align*}
		& \text{Cu} + 2\text{H}_2\text{SO}_4 (\text{конц.}) \rightarrow \text{CuSO}_4 + \text{SO}_2 \uparrow + 2\text{H}_2\text{O} \\
		& 2\text{Al} + 6\text{H}_2\text{SO}_4 (\text{конц.}) \rightarrow \text{Al}_2(\text{SO}_4)_3 + 3\text{SO}_2 \uparrow + 6\text{H}_2\text{O} \\
		& \text{или при сильном нагреве:} \\
		& 2\text{Al} + 4\text{H}_2\text{SO}_4 \rightarrow \text{Al}_2(\text{SO}_4)_3 + \text{S} \downarrow + 4\text{H}_2\text{O} \\
		& 4\text{Al} + 15\text{H}_2\text{SO}_4 \rightarrow 2\text{Al}_2(\text{SO}_4)_3 + 3\text{H}_2\text{S} \uparrow + 12\text{H}_2\text{O}
	\end{align*}
	
	\subsubsection*{Объяснение}
		
	\quad В концентрированной \(H_2SO_4\) окислителем является сера в степени окисления +6 (в составе \(SO_4^{2-}\)). Продукты восстановления зависят от активности металла и условий (концентрация, температура).
	
	\subsubsection*{Опыт 3. Взаимодействие меди с азотной кислотой различной концентрации}
	
	\subsubsection*{Ход работы}
	\begin{enumerate}
		\item В пробирку налито 1–2 мл концентрированной \(HNO_3\) (\(\omega = 63\%\)). Опущены 1–2 стружки меди. Наблюдается интенсивное выделение бурого газа \(NO_2\).
		\item В ту же пробирку добавлена вода до половины объема. Цвет выделяющегося газа изменяется на бесцветный/светло-бурый (\(NO\)).
	\end{enumerate}
	
	\subsubsection*{Уравнения реакций}
	\begin{align*}
		& \text{Cu} + 4\text{HNO}_3 (\text{конц.}) \rightarrow \text{Cu}(\text{NO}_3)_2 + 2\text{NO}_2 \uparrow + 2\text{H}_2\text{O} \\
		& 3\text{Cu} + 8\text{HNO}_3 (\text{разб.}) \rightarrow 3\text{Cu}(\text{NO}_3)_2 + 2\text{NO} \uparrow + 4\text{H}_2\text{O}
	\end{align*}
	
	\subsubsection*{Объяснение}
	
	\quad В азотной кислоте окислителем является азот в степени окисления +5 (в составе \(NO_3^-\)). Продукты восстановления зависят от концентрации кислоты:
	\begin{itemize}
		\item Концентрированная: преимущественно \(NO_2\).
		\item Разбавленная: преимущественно \(NO\) (при дальнейшем разбавлении возможны \(N_2O\), \(NH_4^+\)).
	\end{itemize}
	
	\subsubsection*{Общий вывод}
	
	\quad В ходе лабораторной работы №8 изучены особенности коррозии металлов в растворах кислот:
	\begin{itemize}
		\item Кислоты с катионом-окислителем (\(H^+\) в HCl и разб. \(H_2SO_4\)) окисляют только активные металлы (\(E^0 < 0\)), выделяя водород. Свинец и медь практически не реагируют.
		\item Кислоты с анионом-окислителем (конц. \(H_2SO_4\), \(HNO_3\)) окисляют как активные, так и малоактивные металлы за счет восстановления \(S^{+6}\) или \(N^{+5}\). Продукты восстановления зависят от концентрации кислоты и активности металла.
		\item Концентрированная серная кислота восстанавливается до \(SO_2\) (возможно до \(S\) или \(H_2S\) при сильном нагревании).
		\item Азотная кислота восстанавливается до \(NO_2\) (конц.) или \(NO\) (разб.), что подтверждает зависимость продуктов реакции от концентрации.
	\end{itemize}
	
	\quad Полученные результаты согласуются с теоретическими представлениями об окислительно-восстановительных процессах в кислых средах и влиянии природы металла и кислоты на коррозионное поведение.
	
\end{document}